% Pursuing Stacks (cotes 134-2 to 134-8): archive cote 134-4, site page 98 (PDF page 99), typescript page 185.
% typescript. Draft compared: ps3.tex:4992-5055.
% AI-transcribed by Claude Opus 5.5 (claude-opus-5-5, adaptive thinking, effort high), Message Batches API, batch ps2 (msgbatch_01NVwr1Yuo8qJf5rykeGBXPK), 2026-11-12; not yet checked by eye.
\scanpage{134-4/98}
\mspage{185}

This terminological digression was of a more essential nature, as a
matter of fact, than merely technical. There is still another correcti-%
on I want to make with terminology introduced earlier, namely \xout{the} with
the name of a ``contractor'' I used for intervals endowed with a suitable
composition law (page 120). The name in itself seems to me quite appro-%
priate, however I have now a notion in reserve which seems to me a lot
more important still, a reinforcement it turns out of the notion of a
strict test category - and which I really would like to call a \uline{contrac-}%
\uline{tor}. I could'nt think of any more appropriate name - thus I better
change the previous terminology - sorry ! - and call those nice inter-%
vals ``multiplicative intervals'' , thus referring to the composition law
as a ``multiplication'' (with left unit and left zero object). The name
which first slipped into the typewriter, when it occured a name was desi-%
rable, was not ``contractor'' by the way but ``intersector'' , as I was
thinking of the examples I had met so far, when composition laws were
defined in terms of intersections and were idempotent. But this does'nt
square too well with the example of contractors $\underline{\mathrm{Hom}}(X,X)$, when X is an
object which has a section - and this example turns out as equally sig-%
nificant.

One last comment is about the ``\hand{\v{}}Cech type'' condition L e) on the
basic localizer $\underline{W}$ , introduced \xout{in the} two days ago (page 171). As \deleted{the}
giving\hand{|}\hand{a} ``crible'' in a category amounts to the same as giving a map
\[ C \to \hand{\Delta}_1 \]
(by taking the inverse image of the source-object \hand{\{}0\hand{\}} of $\hand{\Delta}_1$ ), and there-%
fore giving two such amounts to a map from C into $\hand{\Delta}_1 x \hand{\Delta}_1$ , we see that
the \xout{condition} situation when C is the union of two cribles is \xout{the} expres-%
sed \xout{by} equivalently by giving a functor from C into the subcategory
\diagram{\[ C_0 = \hand{\Biggl(} \begin{tikzcd}[row sep=small] & b \\ a \arrow[ur] \arrow[dr] & \\ & c \end{tikzcd} \hand{\Biggr)} \]}
of \hand{$\Delta_1 x \Delta_1$}(dual to the barycentric subdivision of $\hand{\Delta}_1$ ). The asphericity
conditions on C',C'' and C' C'' then jus\hand{t} mean that \uline{this functor is} W-\uline{asphe-}%
\uline{ric}, which by the localization axiom\deleted{s} L 4) implies that the functor it-%
self is a weak equivalence. Thus(by saturation\hand{)}the conclusion that C
should be \inserted{W-}aspheric, just amount\hand{s} to the following condition, which is
in a way the ``universal'' special case when C=C$_0$ , and C',C'' are the
two copies of $\hand{\Delta}_1$ contained in C$_0$ :

L e') The category C$_0$ above is $\underline{W}$-aspheric.

\margin{98}
